O modelo parte dos coeficientes técnicos $A$, da inversa de Leontief $L=(I-A)^{-1}$, da demanda final $y$ e das intensidades $\varphi$. A atribuição de emissões é
\begin{equation}
F=\operatorname{diag}(\varphi)L\operatorname{diag}(y),\qquad c=F\mathbf1,\quad u=F^\top\mathbf1.
\label{eq:artigo-F}
\end{equation}
Cada $F_{ij}$ representa emissões em $i$ atribuídas ao produto final de $j$, em Gg de CO\textsubscript{2}, incorporando requisitos diretos e indiretos. As contas $c$ e $u$ somam as mesmas emissões por origem e destino. O exemplo do Apêndice~\ref{ap:metodo} ilustra essa conservação. As arestas representam atribuições contábeis; concatená-las não reconstrói etapas produtivas.

Para comparar presença como fornecedor e magnitude ambiental, calcula-se $q_i=\sum_{j\ne i}L_{ij}$. Essa soma adimensional dá peso unitário à demanda de cada outro produto final, exclui a diagonal inteira e independe das intensidades ambientais. Não mede quantidade de clientes ou insubstituibilidade. A comparação com as participações em $c$ e $u$ distingue requerimentos padronizados de volumes associados à demanda observada.

A rede dirigida tem 67 nós e pesos $W_{ij}=F_{ij}$ para $i\ne j$, com diagonal zero. As forças externas de saída e entrada são $s_i^{out}=c_i-F_{ii}$ e $s_i^{in}=u_i-F_{ii}$. Sua diferença corresponde ao saldo $u_i-c_i$. A concentração é medida pelo número efetivo $D_2=1/\sum p^2$, em que $p$ representa as participações das contrapartes em cada força. Valores maiores indicam distribuição mais equilibrada; perfis nulos permanecem indefinidos.

O filtro de disparidade \parencite{serrano2009} seleciona relações cuja participação se destaca sob um modelo nulo de repartição aleatória da força. Aplica-se uma vez a $W$, retendo a ligação se o teste na origem ou no destino passa em $\alpha=0{,}05$, com sensibilidade a 0,01 e 0,10. O backbone preserva pesos originais; a cobertura mede a fração de peso retida globalmente e por perfil. As contas e a concentração permanecem calculadas antes do filtro.

